\section{Case studies}
\label{sec:casestudies}

Three case studies show the mechanisms compounding on real workloads: a
fluid-dynamics pipeline assembled almost entirely from open oracles and
compiled classical code (\secref{sec:qfvm}), a symbolic-PDE linearization
whose intermediate artifacts are plans rather than circuits
(\secref{sec:qham-case}), and the reproducible gallery/catalog that pins
behavior across backends (\secref{sec:gallery}).

\subsection{QFVM: compressible Euler equations via reversible flux arithmetic}
\label{sec:qfvm}

Quantum algorithms for fluid dynamics have an active literature~\cite{meng2023fluid,bharadwaj2025fluid}; the quantum finite-volume method (QFVM) targets the one-dimensional
compressible Euler equations by discretizing in a finite-volume form whose
face fluxes are computed from \emph{raw conserved fields} $(\rho, \rho v, E)$
through the frozen-Roe Riemann solver: Roe averages, an entropy fix, and the
eigen-decomposition $R\,|\Lambda|\,R^{-1}$ of the flux Jacobian.
\figref{fig:qfvm} shows the assembled pipeline; every stage is a language
mechanism in action:

\begin{itemize}
\item \textbf{Coherent data access.}  The flow field lives in a QRAM as an
  XOR database of conserved variables---runtime data, locally patchable
  between runs, never embedded in the program (R2, R5).
\item \textbf{Classical physics as compiled reversible code.}  The Roe flux
  formulas---ordinary Python: square roots, divisions, absolute values,
  branching entropy fixes---are compiled by the math-function compiler
  (\secref{sec:mathfunc}) into reversible fixed-point modules with
  propagated status flags.  No matrix-entry lookup table is used: the flux
  is \emph{computed} from the fields, exactly as a classical Roe solver
  computes it.
\item \textbf{Sparse-access assembly.}  The implicit step's operator is
  exposed as a $9$-sparse location/entry oracle pair over the conserved
  state, assembled through a Hermitian dilation of the rectangular flux
  operator plus padding identity---the sparse-access paradigm of
  Table~\ref{tab:paradigms}.
\item \textbf{Injector position.}  The linear solve is performed by
  \emph{any} QLSS protocol under an explicit \texttt{SpectralPromise}: the
  Costa walk kernel and the CKS Chebyshev kernel are interchangeable on the
  same problem object, and the contract machinery rejects inputs whose
  promises are missing.
\end{itemize}

\begin{figure}[t]
  \centering
  \begin{tikzpicture}[
      font=\small,
      stage/.style={draw=black!60, rounded corners=2pt, fill=black!3, align=center, inner sep=4pt, text width=3.0cm},
      mech/.style={font=\scriptsize\itshape, text=black!55, align=center},
      arr/.style={-{Stealth[length=2.2mm]}, black!65},
    ]
    \node[stage, fill=black!8] (qram) at (0, 0)
      {\textbf{QRAM flow fields}\\[-2pt]\scriptsize $\rho,\ \rho v,\ E$\\[-2pt]\scriptsize XOR database, runtime tables};
    \node[stage] (roe) at (5.15, 0)
      {\textbf{Roe flux arithmetic}\\[-2pt]\scriptsize reversible fixed point\\[-2pt]\scriptsize status flags};
    \node[stage] (sparse) at (10.3, 0)
      {\textbf{sparse access}\\[-2pt]\scriptsize location/entry oracles\\[-2pt]\scriptsize Hermitian dilation};
    \node[stage] (qlss) at (10.3, -2.8)
      {\textbf{QLSS kernel}\\[-2pt]\scriptsize Costa walk /\\[-2pt]\scriptsize CKS Chebyshev};
    \node[stage] (read) at (5.15, -2.8)
      {\textbf{norm recovery \& readout}\\[-2pt]\scriptsize joint matrix probe\\[-2pt]\scriptsize host statistics};
    \node[stage] (prom) at (0, -2.8)
      {\textbf{spectral promise}\\[-2pt]\scriptsize $\|A\|$, $\sigma_{\min}$\\[-2pt]\scriptsize evidence strings};

    \draw[arr] (qram) -- node[above, mech] {compiled\\ Python\,$\rightarrow$\,MIR} (roe);
    \draw[arr] (roe) -- node[above, mech] {assembly,\\ dilation} (sparse);
    \draw[arr] (sparse) -- (qlss);
    \draw[arr] (qlss) -- (read);
    \draw[arr] (prom.south) -- ++(0, -0.55) -| (qlss.south);
  \end{tikzpicture}
  \caption{The QFVM pipeline for the one-dimensional compressible Euler
  equations.  Flow fields are served coherently from a QRAM; Roe flux
  formulas written as pure Python are compiled to reversible arithmetic;
  the implicit-step operator is exposed as a sparse-access oracle pair; and
  the linear solve is an \emph{injected} QLSS kernel operating under an
  explicit spectral promise.  Its physical correctness is tracked as
  pending---see the verification matrix of \secref{sec:implementation}.}
  \label{fig:qfvm}
\end{figure}

The pipeline demonstrates the claim that motivated this work: nothing in it
is a monolithic circuit.  Each stage can be bound independently (gate
tables vs.\ QRAM; either QLSS kernel; different fixed-point formats), and
the whole is a module graph of reviewable size.  Its physics-level
correctness is, at the time of writing, an explicitly tracked pending item
(\secref{sec:implementation}).

\subsection{QHAM: from symbolic PDEs to solver-ready input models}
\label{sec:qham-case}

The QHAM path begins at \emph{symbolic} polynomial PDEs---equations are
entered as expressions, not discretized operators---and proceeds: symbolic
equations $\rightarrow$ homotopy-analysis plan (order $m$, parameter
$\eta$) $\rightarrow$ lazy row-rule linearization plan $\rightarrow$
stencil-structured port encodings $\rightarrow$ generator block-encoding
plus norm-preserving lifted initial state $\rightarrow$ any QODE solver.
Two tiers of openness are available to the application: bind the structured
stencil ports, or keep the lifted generator's entire block-encoding open and
carry only the plan in attributes---the open program is then a
\emph{problem specification artifact} that can be bound by whoever owns the
physics (R1).  The general-forced-Burgers example in the repository runs
the full path at order $m=2$, hands the result to a Schr\"odingerization
solver, and exhibits the \texttt{dissipative\_shift} adaptation used when
the downstream family is LCHS or CBMD instead.  Tests verify the
linearization algebra itself (chain-rule identity, stencil-row equality,
and reduction to the plain method at $m=1$)---a division of labor the
evidence ladder of \secref{sec:implementation} makes explicit.

\subsection{The input-model axis}
\label{sec:input-models}

Every solver family of \secref{sec:ode} varies along a second axis
orthogonal to the choice of method: the \emph{input model} through which
problem data enter the program.  The same mathematical instance---a
forced Burgers equation for QHAM, its unforced sibling for Carleman, a
given generator or a discretized heat equation for LCHS and CBMD---can be
presented as structured gate stencils, as a spectral embedding (an exact
Pauli-LCU decomposition of the port or generator matrices), or through
QRAM-served data paths: coefficient angle databases, prepare tables,
sparse-access position/entry tables, and QRAM-loaded initial states.  The
effect of the choice is visible in the artifacts themselves rather than
in any solver code: which oracle slots remain open in the exported
program, which QRAM resources appear after binding, and whether the data
are exact gate truth tables or angle-quantized words at a declared
precision.  Appendix~\ref{app:input-models} runs all four families
through this matrix on common instances and records the resulting
open-slot and resource signatures (\tabref{tab:input-models}).

Two consequences are worth noting here.  First, the QRAM variants are
not special-purpose programs: the coefficient-angle database is produced
by one new adapter, \texttt{qram\_coefficient\_encoding}, plugged into
the existing \texttt{coefficient\_encoder} hook of the structured
finite-difference bindings, and the sparse-access grid stencil is the
same \texttt{SparseA} paradigm of \secref{sec:oracles} promoted to a
\emph{discretization}---so a QRAM-served grid flows through
\texttt{make\_qpde} into any QODE family unchanged.  Second, input models
compose across families: the open coefficient-database QHAM input, after
\texttt{dissipative\_shift}, is solved by CBMD with no change to its
bindings, and the same open program closes once against a gate database
and once against a QRAM database with only the runtime memory table
changed (R2).

\subsection{The gallery and the reproducible catalog}
\label{sec:gallery}

Two artifacts anchor day-to-day verification:

\begin{itemize}
\item \textbf{The algorithm gallery} (22 cases) spans the base library:
  query algorithms, Fourier arithmetic, search and amplification,
  estimation (phase, amplitude, counting, Hadamard, swap tests), variational
  circuit generation, quantum walks, modular arithmetic and order finding,
  repetition codes, and Trotter evolution.  Every case is exported to RIR,
  OriginIR-ext, and the strict gate set, and---in the native
  configuration---its full amplitude vector is cross-checked between the
  reference simulator and \emph{both} real backends (PySparQ and the
  OriginIR/uniqc stack) to $10^{-10}$.  This is the circuit-witness level of
  the evidence ladder: it certifies that what the language generated is what
  the backends executed.
\item \textbf{The applications catalog} (32 reproducible cases) covers the
  scientific-computing stack: sparse and database input models in gate
  \emph{and} QRAM forms, both QLSS kernels, the ODE families, QHAM and QFVM
  entries, QSVT transforms, amplitude amplification, block-encoding
  algebra, arithmetic, batched and banked QRAM, register views, and
  measurement plans.  Each catalog entry is an \emph{open} program with a
  declared binding plan and runtime memory tables; integration tests parse
  every exported artifact with the real OriginIR parser.
\end{itemize}

Together they pin the meaning of ``the same program'' across realizations
(R2) and backends, and they delimit honestly what is verified: structural
validity and circuit-level agreement everywhere; numerical certification of
the scientific solvers tracked separately, as the next section details.
